\documentclass[10pt,a4paper]{amsart}
\usepackage{amsmath,amssymb,amsthm,graphicx}
\usepackage{verbatim}
\newenvironment{reference}{\comment}{\endcomment}
%\newenvironment{reference}{}{}
\newenvironment{slogan}{\comment}{\endcomment}
\newenvironment{history}{\comment}{\endcomment}

% For commutative diagrams we use Xy-pic
\usepackage[all]{xy}

% We use 2cell for 2-commutative diagrams.
\xyoption{2cell}
% Theorem environments.
%
\theoremstyle{plain}
\newtheorem{theorem}[subsection]{Theorem}
\newtheorem{proposition}[subsection]{Proposition}
\newtheorem{lemma}[subsection]{Lemma}

\theoremstyle{definition}
\newtheorem{definition}[subsection]{Definition}
\newtheorem{example}[subsection]{Example}
\newtheorem{exercise}[subsection]{Exercise}
\newtheorem{situation}[subsection]{Situation}

\theoremstyle{remark}
\newtheorem{remark}[subsection]{Remark}
\newtheorem{remarks}[subsection]{Remarks}

\numberwithin{equation}{subsection}

% Macros
%
\def\lim{\mathop{\mathrm{lim}}\nolimits}
\def\colim{\mathop{\mathrm{colim}}\nolimits}
\def\Spec{\mathop{\mathrm{Spec}}}
\def\Hom{\mathop{\mathrm{Hom}}\nolimits}
\def\Ext{\mathop{\mathrm{Ext}}\nolimits}
\def\SheafHom{\mathop{\mathcal{H}\!\mathit{om}}\nolimits}
\def\SheafExt{\mathop{\mathcal{E}\!\mathit{xt}}\nolimits}
\def\Sch{\mathit{Sch}}
\def\Mor{\mathop{\mathrm{Mor}}\nolimits}
\def\Ob{\mathop{\mathrm{Ob}}\nolimits}
\def\Sh{\mathop{\mathit{Sh}}\nolimits}
\def\NL{\mathop{N\!L}\nolimits}
\def\CH{\mathop{\mathrm{CH}}\nolimits}
\def\proetale{{pro\text{-}\acute{e}tale}}
\def\etale{{\acute{e}tale}}
\def\QCoh{\mathit{QCoh}}
\def\Ker{\mathop{\mathrm{Ker}}}
\def\Im{\mathop{\mathrm{Im}}}
\def\Coker{\mathop{\mathrm{Coker}}}
\def\Coim{\mathop{\mathrm{Coim}}}

% Boxtimes
%
\DeclareMathSymbol{\boxtimes}{\mathbin}{AMSa}{"02}

%
% Macros for moduli stacks/spaces
%
\def\QCohstack{\mathcal{QC}\!\mathit{oh}}
\def\Cohstack{\mathcal{C}\!\mathit{oh}}
\def\Spacesstack{\mathcal{S}\!\mathit{paces}}
\def\Quotfunctor{\mathrm{Quot}}
\def\Hilbfunctor{\mathrm{Hilb}}
\def\Curvesstack{\mathcal{C}\!\mathit{urves}}
\def\Polarizedstack{\mathcal{P}\!\mathit{olarized}}
\def\Complexesstack{\mathcal{C}\!\mathit{omplexes}}
% \Pic is the operator that assigns to X its picard group, usage \Pic(X)
% \Picardstack_{X/B} denotes the Picard stack of X over B
% \Picardfunctor_{X/B} denotes the Picard functor of X over B
\def\Pic{\mathop{\mathrm{Pic}}\nolimits}
\def\Picardstack{\mathcal{P}\!\mathit{ic}}
\def\Picardfunctor{\mathrm{Pic}}
\def\Deformationcategory{\mathcal{D}\!\mathit{ef}}
\usepackage{cleveref}
\crefname{originallemma}{original lemma}{original lemmas}
\Crefname{originallemma}{Original lemma}{Original lemmas}
\makeatletter
\newlabel{lemma-flat-eq}{{00HK}{source}{}{sourcepointer-1}{}}
\newlabel{lemma-flat-eq@cref}{{[originallemma][0][]00HK}{[1][0][]source}{}{}{}}
\makeatother
\crefname{originallemma}{original lemma}{original lemmas}
\Crefname{originallemma}{Original lemma}{Original lemmas}
\makeatletter
\newlabel{lemma-local-criterion-flatness}{{00MK}{source}{}{sourcepointer-2}{}}
\newlabel{lemma-local-criterion-flatness@cref}{{[originallemma][0][]00MK}{[1][0][]source}{}{}{}}
\makeatother
\crefname{originallemma}{original lemma}{original lemmas}
\Crefname{originallemma}{Original lemma}{Original lemmas}
\makeatletter
\newlabel{lemma-what-does-it-mean}{{051C}{source}{}{sourcepointer-3}{}}
\newlabel{lemma-what-does-it-mean@cref}{{[originallemma][0][]051C}{[1][0][]source}{}{}{}}
\makeatother
\crefname{originallemma}{original lemma}{original lemmas}
\Crefname{originallemma}{Original lemma}{Original lemmas}
\makeatletter
\newlabel{remark-Tor-ring-mod-ideal}{{00M6}{source}{}{sourcepointer-4}{}}
\newlabel{remark-Tor-ring-mod-ideal@cref}{{[originallemma][0][]00M6}{[1][0][]source}{}{}{}}
\makeatother
\title{A variant of the local criterion for flatness via Tor one and the quotient by an ideal}
\author{The Stacks Project Authors}
\date{stacks-project:algebra.tex:local-snapshot-20261012; GNU FDL 1.2 or later}
\begin{document}
\maketitle
\noindent\textit{Editorial note.} All delivered mathematical text is the author's own, copied from the pinned Stacks Project chapter file algebra.tex; no statement or proof step was rewritten, reordered or removed.\par\medskip
\noindent\textit{Editorial note.} Preamble provenance, as required for this source: the frozen snapshot stages the chapter's own preamble file beside it (preamble.tex, 217 lines, 8116 bytes, the same shared Stacks-wide preamble the published chow and etale-cohomology cards use); the first line of the pinned chapter file is \textbackslash{}input\{preamble\}, so no preamble text lives inside the chapter itself. The card takes that file's lines 8-18 (the verbatim package and the three \textbackslash{}comment-based environments reference/slogan/history at lines 9-12, section 13.2, together with the Xy-pic setup) and its lines 154 to the runtime end of the file (the native theorem declarations and mathematical macros). Not one character was changed, no package was added, no macro was defined, and the excerpt serves only as preamble, never as a span coordinate.\par\medskip
\noindent\textit{Editorial note.} The same-chapter pointers are located at their \textbackslash{}label lines inside the pinned chapter file; the cross-chapter pointers are located in the project's own tag list staged in the snapshot (sources/stacks-chow/tags.txt), at the verbatim lines carrying their canonical tags.\par\medskip
\noindent\textit{Editorial note.} The counted proof refers to same-chapter statements by label; they are recorded as navigation stubs pointed at their \textbackslash{}label lines, and the arguments stay inside the counted span.\par\medskip
\noindent\textit{Editorial note.} The counted spans contain no \textbackslash{}cite, so the card carries no bibliography.\par\medskip
\noindent\textit{Allowed standard prerequisites (editorial).}\begin{itemize}
\item Vanishing of Tor one against the residue field implies flatness. Invoked in the first proof inside the span (lines 23976-23981). Reduces the claim to a Tor computation against the residue field.
\item Tor one against a quotient as a kernel. Used in the second proof inside the span (lines 23990-24014). Reinterprets the hypothesis as injectivity of a multiplication map.
\end{itemize}\medskip
\ifdefined\hypertarget\hypertarget{sourcepointer-1}{}\fi
\noindent\textit{Source reference.} Original lemma 00HK: sources/stacks-algebra/pinned-original-source.tex, source lines 9212--9212. Navigation pointer for \textbackslash{}ref\{lemma-flat-eq\} used inside the counted span (line 23990); the argument that consumes it is retained in the span.\par\medskip
\ifdefined\hypertarget\hypertarget{sourcepointer-2}{}\fi
\noindent\textit{Source reference.} Original lemma 00MK: sources/stacks-algebra/pinned-original-source.tex, source lines 23800--23800. Navigation pointer for \textbackslash{}ref\{lemma-local-criterion-flatness\} used inside the counted span (line 23981); the argument that consumes it is retained in the span.\par\medskip
\ifdefined\hypertarget\hypertarget{sourcepointer-3}{}\fi
\noindent\textit{Source reference.} Original lemma 051C: sources/stacks-algebra/pinned-original-source.tex, source lines 23868--23868. Navigation pointer for \textbackslash{}ref\{lemma-what-does-it-mean\} used inside the counted span (line 23977); the argument that consumes it is retained in the span.\par\medskip
\ifdefined\hypertarget\hypertarget{sourcepointer-4}{}\fi
\noindent\textit{Source reference.} Original lemma 00M6: sources/stacks-algebra/pinned-original-source.tex, source lines 18907--18907. Navigation pointer for \textbackslash{}ref\{remark-Tor-ring-mod-ideal\} used inside the counted span (line 24014); the argument that consumes it is retained in the span.\par\medskip
\begin{lemma}[Variant of the local criterion]
\label{lemma-variant-local-criterion-flatness}
Let $R \to S$ be a local homomorphism of Noetherian
local rings. Let $I \not = R$ be an ideal in $R$.
Let $M$ be a finite $S$-module. If $\text{Tor}_1^R(M, R/I) = 0$
and $M/IM$ is flat over $R/I$, then $M$ is flat over $R$.
\end{lemma}

\begin{proof}
First proof: By
Lemma \ref{lemma-what-does-it-mean}
we see that $\text{Tor}_1^R(\kappa, M)$ is zero where $\kappa$
is the residue field of $R$. Hence we see that $M$
is flat over $R$ by
Lemma \ref{lemma-local-criterion-flatness}.

\medskip\noindent
Second proof: Let $\mathfrak m$ be the maximal ideal of $R$.
We will show that $\mathfrak m \otimes_R M \to M$ is injective,
and then apply
Lemma \ref{lemma-local-criterion-flatness}.
Suppose that $\sum f_i \otimes x_i \in \mathfrak m \otimes_R M$
and that $\sum f_i x_i = 0$ in $M$. By the equational criterion
for flatness Lemma \ref{lemma-flat-eq} applied to $M/IM$
over $R/I$ we see there exist $\overline{a}_{ij} \in R/I$
and $\overline{y}_j \in M/IM$ such that
$x_i \bmod IM = \sum_j \overline{a}_{ij} \overline{y}_j $
and $0 = \sum_i (f_i \bmod I) \overline{a}_{ij}$.
Let $a_{ij} \in R$ be a lift of $\overline{a}_{ij}$ and
similarly let $y_j \in M$ be a lift of $\overline{y}_j$.
Then we see that
\begin{eqnarray*}
\sum f_i \otimes x_i
& = &
\sum f_i \otimes x_i +
\sum f_ia_{ij} \otimes y_j -
\sum f_i \otimes a_{ij} y_j
\\
& = &
\sum f_i \otimes (x_i - \sum a_{ij} y_j) +
\sum (\sum f_i a_{ij}) \otimes y_j
\end{eqnarray*}
Since $x_i - \sum a_{ij} y_j \in IM$ and
$\sum f_i a_{ij} \in I$ we see that there exists
an element in $I \otimes_R M$ which maps to our given
element $\sum f_i \otimes x_i$ in $\mathfrak m \otimes_R M$.
But $I \otimes_R M \to M$ is injective by assumption (see
Remark \ref{remark-Tor-ring-mod-ideal}) and we win.
\end{proof}

\end{document}
